\section{Erd\H{o}s Problem \#778: independent continuation}
\noindent Date: 2026-09-23. Research attempt; no claim of a new published result.
\subsection{1) TARGET AND SOURCE AUDIT}
In the complete-graph edge claiming games, Alice claims one red edge per move and Bob claims blue edges, with ties awarded to Bob. Compare their largest cliques, also considering the maximum-degree (star) variant and Bob bias two.

All supplied versions considered: \texttt{778.tex}.
The source imports the symmetric outcomes through $n=8$ and retains biased Clique$^{(1,2)}(n)$ only through $n=6$. I independently implement exact minimax and extend the checked biased range to $n=7$. This is a finite computational result, with no claim of a general strategy.
\subsection{2) WORK}
\paragraph{Exact recursion.}
Index the edges of $K_n$ by bit positions. A state is $(R,B,t)$, the disjoint red and blue masks and the player to move. At terminal states return true precisely when Alice has the strictly larger score. For a clique score, test all vertex subsets for completeness; for a star score, take maximum degree. The recursion is
\[
W(R,B,A)=\bigvee_{e\ \mathrm{free}}W(R\cup\{e\},B,B\mathrm{ob}),
\]
\[
W(R,B,B\mathrm{ob})=\bigwedge_{S\subseteq E_{\rm free},\ |S|=\min(b,|E_{\rm free}|)}W(R,B\cup S,A),
\]
where $b$ is Bob's bias. An induction on the number of unclaimed edges proves the recursion exactly represents optimal play. Bob's consecutive selections form an unordered batch, since Alice does not move between them.

The executable \texttt{verification/batch\_2\_games.py} evaluates this recursion with memoization and short-circuit logical operations. Actual results: Bob wins both clique biases for $3\le n\le5$, Alice wins the star game at $n=3$, and Bob wins it at $n=4,5$. For bias two, Bob also wins at $n=6$ and $n=7$. The $n=6$ run visits 5,947 memoized states; the $n=7$ run visits 199,911. Thus the source\textquotesingle s independently checked biased range extends to
\[
\boxed{\text{Bob wins Clique}^{(1,2)}(n)\quad(3\le n\le7).}
\]
This is a reproducible finite certificate by exhaustive recursive evaluation, not an assertion that all ternary colorings were explicitly visited. Unvisited branches are skipped only when the Boolean recurrence already determines the value.
\subsection{3) ADVERSARIAL VERIFICATION}
Ties go to Bob, Alice moves first, and a final incomplete Bob move takes every remaining edge. These conventions are implemented explicitly. There is no symmetry reduction or heuristic pruning. A single implementation is evidence with an inspectable exhaustive algorithm, not a formally verified game proof; the source\textquotesingle s symmetric $n=7,8$ computations were not reproduced here.
\subsection{4) FINAL}
\textbf{UNRESOLVED.}
No winning strategy for arbitrary $n$ is obtained for either bias or for the star variant. The new bounded computation cannot justify extrapolation.
\subsection{5) SOURCES CHECKED}
Full source read. Checked Cambie--Provoost, \texttt{https://arxiv.org/html/2505.03497v2}, Proposition 9, on 2026-09-23, confirming the symmetric $3\le n\le8$ range. The biased $n=7$ result above comes from the local executable computation.
\subsection{6) COMPLETION ESTIMATE}
This is a conservative subjective measure of progress toward the original question, not a probability of correctness or a claim that the remaining work is routine.
\textbf{COMPLETION: 12\%}
